\chapter{逻辑、集合与映射}\label{ch:b1:logic}

到目前为止，证明都是在一种非形式化、但诚实的
“证明”概念下进行的。本章是大学数学的第一章，
它将游戏规则讲清楚：什么是数学命题，
联结词与量词如何组合命题，证明中哪些推理是合法的——
然后在此基础上，
建立数学的两门通用语言：\omterm{def:b1:logic:sets}{集合}与\omterm{def:b1:logic:map}{映射}。

\section{命题与联结词}

\begin{definition}[命题与联结词]\label{def:b1:logic:statement}
称一个句子为\emph{命题}\index{命题}（亦称\emph{命题}）
，它必为真（T）或为假（F）——二者恰居其一。由命题 $P$ 与
$Q$ 可以构造出下列命题：
\begin{itemize}
  \item \emph{否定} $\lnot P$（“非 $P$”），
        仅当 $P$ 为假时为真；
  \item \emph{合取} $P \land Q$（“$P$ 且 $Q$”），
        仅当两者都为真时为真；
  \item \emph{析取} $P \lor Q$（“$P$ 或 $Q$”），
        仅当至少一者为真时为真（此处的“或”是可兼的）；
  \item \emph{蕴含} $P \implies Q$，仅当 $P$
        为真而 $Q$ 为假时它为假；
  \item \emph{等价} $P \iff Q$，仅当 $P$
        与 $Q$ 真值相同时为真。
\end{itemize}
\end{definition}

\begin{remark}
$P \implies Q$ 的真值表值得注意：当 $P$ 为假时，
$P \implies Q$ 无论 $Q$ 为何都为\emph{真}。“若 $2 < 1$，则
“$0 = 5$”是一个真的蕴含。蕴含对
其前提不成立时发生的事情不作任何断言。
\end{remark}

\begin{proposition}[命题演算的规则]\label{prop:b1:logic:rules}
对任意命题 $P$、$Q$、$R$：
\begin{enumerate}
  \item $\lnot(\lnot P) \iff P$；
  \item 德摩根律\index{德摩根律}：
        $\lnot(P \land Q) \iff (\lnot P) \lor (\lnot Q)$ 与
        $\lnot(P \lor Q) \iff (\lnot P) \land (\lnot Q)$；
  \item $(P \implies Q) \iff \bigl((\lnot P) \lor Q\bigr)$，因而
        $\lnot(P \implies Q) \iff P \land (\lnot Q)$；
  \item 逆否律\index{逆否律}：
        $(P \implies Q) \iff \bigl((\lnot Q) \implies (\lnot P)\bigr)$；
  \item $(P \iff Q) \iff \bigl((P \implies Q) \land (Q \implies
        P)\bigr)$；
  \item 分配律：$P \land (Q \lor R) \iff (P \land Q) \lor
        (P \land R)$ 与 $P \lor (Q \land R) \iff (P \lor Q) \land
        (P \lor R)$。
\end{enumerate}
\end{proposition}

\begin{proof}
每个等价式都通过比较真值表来验证：两个
由 $P$、$Q$、$R$ 构成的复合命题等价，当且仅当它们
在每一种（四种或八种）情形下真值相同。下面完整列出
第一条德摩根律的真值表：
\begin{center}
\begin{tabular}{cc|c|c|cc|c}
$P$ & $Q$ & $P \land Q$ & $\lnot(P \land Q)$ & $\lnot P$ & $\lnot Q$ &
$(\lnot P) \lor (\lnot Q)$ \\
\hline
T & T & T & F & F & F & F \\
T & F & F & T & F & T & T \\
F & T & F & T & T & F & T \\
F & F & F & T & T & T & T \\
\end{tabular}
\end{center}
第 $4$ 列与第 $7$ 列相同，由此证得该律。对于逆否律
用文字论证更为简便：$P \implies Q$ 仅在
（$P$ 真、$Q$ 假）时为假，而 $(\lnot Q) \implies (\lnot P)$ 仅在
（$\lnot Q$ 真、$\lnot P$ 假）时为假，即（$Q$
假、$P$ 真）——同一个单一情形，因此两个蕴含的
真值表相同。其余规则同法验证；注意到
（3）把每个蕴含化为析取，因此由（2）
机械地得出否定规则 $\lnot(P \implies Q) \iff P
\land (\lnot Q)$：要推翻一个蕴含，必须给出一种情形，
使前提成立而结论不成立。
\end{proof}

\section{量词}

\begin{definition}[量词]\label{def:b1:logic:quantifiers}
设 $P(x)$ 为元素 $x$ 的一个性质，该元素属于\omterm{def:b1:logic:sets}{集合} $E$。
\begin{itemize}
  \item $\forall x \in E,\ P(x)$（“对所有 $x$ 属于 $E$，$P(x)$”）为真，
        当 $E$ 中每个元素都满足 $P$ 时为真；
  \item $\exists x \in E,\ P(x)$（“存在 $x$ 属于 $E$，使得
        $P(x)$”）为真，当 $E$ 中至少有一个元素满足
        $P$。
\end{itemize}
用 $\exists!$ 表示“存在唯一”。
\end{definition}

\begin{proposition}[量词的否定]\label{prop:b1:logic:negquant}
\[
\lnot\bigl(\forall x \in E,\ P(x)\bigr) \iff
\exists x \in E,\ \lnot P(x),
\qquad
\lnot\bigl(\exists x \in E,\ P(x)\bigr) \iff
\forall x \in E,\ \lnot P(x).
\]
\end{proposition}

\begin{proof}
下面分两个方向证明第一个等价式；第二个
与之对称。若 $\forall x \in E,\ P(x)$ 不成立，则并非每个
元素都满足 $P$：\omterm{def:b1:logic:sets}{集合} $A = \{x \in E : \lnot P(x)\}$
非空，它的任一元素都见证了 $\exists x \in
E,\ \lnot P(x)$。反之，若某个 $x_0 \in E$ 满足 $\lnot
P(x_0)$，则 $x_0$ 即为反例，全称命题
不成立。对于第二条规则，“没有 $x$ 满足 $P$”表示\omterm{def:b1:logic:sets}{集合}
$\{x : P(x)\}$ 为空，即每个 $x$ 都属于它的补集
$A$。把这两条规则依次作用于嵌套的量词前缀，
便得到\cref{ex:b1:logic:limit}所给的机械步骤：
否定自左向右进行，把每个 $\forall$ 换为
$\exists$，每个 $\exists$ 换为 $\forall$，最后
对最内层的谓词取否定。
\end{proof}

\begin{example}[日常数学语句的否定]\label{ex:b1:logic:negpractice}
设 $f \colon \R \to \R$。句子“$f$ 是增函数”写作
\[
\forall x \in \R,\ \forall y \in \R,\quad
x \leq y \implies f(x) \leq f(y) ,
\]
其否定由\cref{prop:b1:logic:negquant}并结合规则
$\lnot(P \implies Q) \iff P \land \lnot Q$，得：
\[
\exists x \in \R,\ \exists y \in \R,\quad
x \leq y \ \text{ 且 }\ f(x) > f(y) :
\]
只需一对见证即可。同理，“$f$ 有界”即
$\exists M \in \R,\ \forall x \in \R,\ \abs{f(x)} \leq M$，其
否定为
\[
\forall M \in \R,\ \exists x \in \R,\quad \abs{f(x)} > M :
\]
\emph{无论}所提出的界是多少，总有某点超过它。要点
在于：正确的否定绝不对整个量化块施加“非”——
它本身是一个新的肯定陈述，其中角色互换：
现在要给出的，正是此前所接收的
见证。
\end{example}

\begin{example}[量词的次序]\label{ex:b1:logic:order}
不同量词的次序是有影响的：
\[
\forall x \in \R,\ \exists y \in \R,\ y > x
\quad\text{为真（取 } y = x+1\text{），}
\]
\[
\exists y \in \R,\ \forall x \in \R,\ y > x
\quad\text{为假（不存在超过所有实数的实数）。}
\]
在第一个命题中，$y$ 可依赖于 $x$；在第二个中，唯一的
$y$ 须对所有 $x$ 成立。另一方面，两个相同的量词
则总可以交换次序。
\end{example}

\begin{example}[用三个量词读懂定义]\label{ex:b1:logic:limit}
句子“数列 $(u_n)$ 收敛于 $\ell$”将
在\cref{ch:b1:seq}中写作
\[
\forall \varepsilon > 0,\ \exists N \in \N,\ \forall n \geq N,\quad
\abs{u_n - \ell} \leq \varepsilon .
\]
其否定由\cref{prop:b1:logic:negquant}连用三次而得，为
\[
\exists \varepsilon > 0,\ \forall N \in \N,\ \exists n \geq N,\quad
\abs{u_n - \ell} > \varepsilon .
\]
能够机械地否定此类语句，无须考虑其含义，
这是一项真正的技能，它把逻辑工作
与数学工作区分开来。
\end{example}

\section{证明技巧}

\begin{method}[标准的证明模式]\label{met:b1:logic:proofs}
要证明\dots
\begin{enumerate}
  \item \emph{直接证明蕴含 $P \implies Q$}：设 $P$，
        由此推出 $Q$；
  \item \emph{用逆否命题证明}：设 $\lnot Q$，推出 $\lnot P$
——由\cref{prop:b1:logic:rules}~(4)可知其有效；
  \item \emph{用反证法证明}：假设命题不成立，
        推出矛盾；
  \item \emph{证明等价}：分别证明两个方向的蕴含（或
        串接已知的等价）；
  \item \emph{全称命题}：取\emph{任意}的 $x$
        属于 $E$（记为“设 $x \in E$”），并证明 $P(x)$；
  \item \emph{存在性命题}：给出一个见证，或
        间接证明其存在；
  \item \emph{用归纳法}：见\cref{thm:b1:logic:induction}。
\end{enumerate}
证明关于一个选定但任意的元素的命题时，
切勿给元素附加额外性质：“设 $x \in \R$”之后
再写“既然 $x > 0$……”，对负的 $x$ 什么也证明不了。
\end{method}

\begin{remark}[证明中的常见误区]\label{rem:b1:logic:pitfalls}
四个经典陷阱，各值得点明一次。
\begin{enumerate}
  \item \emph{以逆命题代替逆否命题。} $Q \implies P$
        \emph{不}等价于 $P \implies Q$；只有 $\lnot Q
        \implies \lnot P$ 才与之等价。“若下雨，则街道湿”
        并不能由街道湿而断定下过雨。
  \item \emph{用单个蕴含证明等价。}一个“当且仅当”断言
        实为两个定理；须说明正在证明的方向，
        两个方向都要证明。$\iff$ 的链式推理只有在
        \emph{每一步}确实可逆时才合法——
        例如两边平方一个方程，便并不可逆。
  \item \emph{逆向证明。}从所希望的结论出发，
        推出一个真命题并不能证明任何东西
        （由 $-1 = 1$ 两边平方，得真命题 $1 = 1$）。
        计算可以\emph{逆向发现}，但它
        必须\emph{正向写出}，或附上明确的
        等价变形。
  \item \emph{固定见证与任意元素。}要证明
        $\exists x,\ P(x)$，可以给出一个巧妙选取的
        $x$；要证明 $\forall x,\ P(x)$，所取的 $x$
        必须保持任意。混淆二者——
        以一个例子检验全称命题——
        是初学者作业中最常见的错误。
\end{enumerate}
\end{remark}

\begin{example}[逆否命题与反证法的应用]\label{ex:b1:logic:sqrt2}
\emph{对 $n \in \N$：若 $n^2$ 为偶数，则 $n$ 为偶数。}用
逆否命题证明：若 $n$ 为奇数，$n = 2k+1$，则
$n^2 = 4k^2 + 4k + 1$ 为奇数。

\emph{$\sqrt 2$ 是无理数。}用反证法：假设
$\sqrt 2 = p/q$，其中 $p, q \in \N^*$，且该分数为既约分数。
于是 $p^2 = 2q^2$ 为偶数，故 $p$ 为偶数（见上文），$p = 2r$；
则 $q^2 = 2r^2$ 为偶数，故 $q$ 为偶数——与分数
既约的假设矛盾。
\end{example}

\begin{theorem}[数学归纳法]\label{thm:b1:logic:induction}
设 $P(n)$ 为整数 $n$ 的性质。若
\begin{enumerate}
  \item $P(0)$ 为真，且
  \item 对所有 $n \in \N$，$P(n) \implies P(n+1)$，
\end{enumerate}
则 $P(n)$ 对所有 $n \in \N$ 为真。

\emph{强归纳法：}将 (2) 换为下式，结论不变：
对所有 $n$，$\bigl(P(0) \land \dots \land P(n)\bigr) \implies
P(n+1)$.
\end{theorem}

\begin{proof}
这是 $\N$ 本身的性质，等价于：\emph{$\N$ 的每个非空
子集必有最小元}（此处视为已知）。事实上，
假设 (1) 与 (2) 成立，令 $A = \{n \in \N : P(n)
\text{ 不成立}\}$。若 $A \neq \emptyset$，则它有最小元 $m$；
由 (1) 知 $m \neq 0$；于是 $m - 1 \notin A$，故 $P(m-1)$ 成立，由 (2)
得 $P(m)$——矛盾。故 $A = \emptyset$。对于强
归纳法，同样论证：$P(0), \dots, P(m-1)$ 均成立，
因为 $m$ 是 $A$ 中的最小元。
\end{proof}

\begin{example}[证明唯一存在性]\label{ex:b1:logic:uniqueexistence}
命题 $\exists!\,x,\ P(x)$ 实为\emph{两个}命题，
需分别证明：存在性（找出或构造某个 $x_0$ 使
$P(x_0)$ 成立）与唯一性（设 $P(x)$ 与 $P(x')$ 成立，推出 $x =
x'$）。例如：\emph{存在唯一的实数 $x$，使得 $x^3 + x =
2$。}存在性：取 $x_0 = 1$ 即可，因为 $1 + 1 = 2$。唯一性：
若 $x^3 + x = x'^3 + x'$，则
\[
0 = (x^3 - x'^3) + (x - x')
= (x - x')\,\bigl(x^2 + xx' + x'^2 + 1\bigr),
\]
且第二个因子为正（它等于 $\bigl(x +
\tfrac{x'}2\bigr)^2 + \tfrac34 x'^2 + 1 \geq 1$），故 $x = x'$。
注意分工：存在性用了一个幸运的猜测，
唯一性则用了对\emph{任意}解成立的代数推理——
两个论证互不替代；找到解之后，忘掉后
半部分是一种常见的诱惑。
\end{example}

\begin{example}\label{ex:b1:logic:induction}
对所有 $n \in \N^*$：$\;\sum_{k=1}^n k = \frac{n(n+1)}{2}$。归纳基础
$n = 1$：两边均为 $1$。归纳步：假设公式对 $n$ 成立，
\[
\sum_{k=1}^{n+1} k = \frac{n(n+1)}{2} + (n+1)
= (n+1)\Bigl(\frac n2 + 1\Bigr) = \frac{(n+1)(n+2)}{2}. \qedhere
\]
\end{example}

\begin{example}[强归纳法的应用]\label{ex:b1:logic:strongind}
\emph{每个整数 $n \geq 2$ 都可写成质数之积}（质数是指
$\geq 2$ 的整数，且其 $\geq 1$ 的因数只有 $1$ 与
自身；质数本身的性质在\cref{ch:b1:arith}中研究）。
普通归纳法在此无能为力：已知 $95 = 5 \times 19$
的分解对 $96$ 毫无帮助。强归纳法恰好适用。归纳
基础：$2$ 是质数，因而是仅含一个因子的质数之积。归纳步：设
$n \geq 2$，并假设每个整数 $m$（$2 \leq m \leq n$）都是
质数之积。若 $n + 1$ 为质数，则结论成立；否则 $n + 1 = ab$
其中 $2 \leq a, b \leq n$；由强归纳假设，$a$ 与 $b$ 都是
质数之积，从而 $n + 1$ 亦然。要点在于：强
归纳法是最合适的工具，每当“$P(n+1)$ 的理由”
位于某个难以预料的较早序号，而非序号 $n$。
\end{example}

\section{集合}

\begin{definition}[集合的运算]\label{def:b1:logic:sets}
将\emph{集合}\index{集合}的概念与属于关系
$x \in E$ 作为原始概念。设 $A, B$ 为全集 $E$ 中的集合：
\begin{itemize}
  \item \emph{包含}：$A \subseteq B$ 指 $\forall x,\ x \in A
        \implies x \in B$；相等 $A = B$ 当 $A \subseteq B$ 且
        $B \subseteq A$；
  \item \emph{并集} $A \cup B$，\emph{交集} $A \cap B$，
        \emph{差集} $A \setminus B = \{x \in A : x \notin B\}$，
        \emph{补集} $\overline{A} = E \setminus A$；
  \item \emph{空集} $\emptyset$，包含于每个集合；
  \item \emph{幂集}\index{幂集} $\mathcal{P}(E)$：由
        $E$ 的全部子集组成的集合；
  \item \emph{笛卡儿积} $E \times F$：由有序对
        $(x, y)$ 组成的集合，其中 $x \in E$，$y \in F$。
\end{itemize}
\end{definition}

\begin{example}[熟悉幂集]\label{ex:b1:logic:powerset}
设 $E = \{a, b\}$：
\[
\mathcal P(E) = \bigl\{\, \emptyset,\ \{a\},\ \{b\},\ \{a, b\}
\,\bigr\},
\]
共四个元素——注意其类型的区别：$a \in E$，但
$\{a\} \in \mathcal P(E)$；命题 $a \in \mathcal P(E)$
与 $\{a\} \subseteq \mathcal P(E)$ 按此写法均为假
（后者要求 $a$ 为 $E$ 的\emph{子集}）。
从空集出发反复求\omterm{def:b1:logic:sets}{幂集}：$\mathcal P(\emptyset) = \{\emptyset\}$
只有一个元素，$\mathcal P(\mathcal P(\emptyset)) = \{\emptyset,
\{\emptyset\}\}$ 有两个元素，下一个有四个——\omterm{def:b1:logic:sets}{集合}之\omterm{def:b1:logic:sets}{集合}
仍是普通\omterm{def:b1:logic:sets}{集合}，这一倍增规律将在\cref{ch:b1:counting}中
得到证实：$\abs{\mathcal P(E)} = 2^{\abs E}$。分清各层次
（$x$、$\{x\}$、$\{\{x\}\}$）是解题的关键之一，
如\cref{exo:b1:logic:11,exo:b1:logic:12}。
\end{example}

\begin{proposition}[集合代数]\label{prop:b1:logic:setalgebra}
设 $A, B, C$ 为 $E$ 的子集：
\begin{enumerate}
  \item $A \cap (B \cup C) = (A \cap B) \cup (A \cap C)$ 且
        $A \cup (B \cap C) = (A \cup B) \cap (A \cup C)$；
  \item 德摩根律：$\overline{A \cup B} = \overline{A} \cap
        \overline{B}$ 且 $\overline{A \cap B} = \overline{A} \cup
        \overline{B}$；
  \item $A \subseteq B \iff \overline{B} \subseteq \overline{A}$。
\end{enumerate}
\end{proposition}

\begin{proof}
每个恒等式都是把\cref{prop:b1:logic:rules}中的规则经由
对应关系（$\in A$ 或否定）$\leftrightarrow$（命题
真或假）翻译得到。例如 $x \in \overline{A \cup B} \iff \lnot(x \in A \lor x \in
B) \iff (x \notin A) \land (x \notin B) \iff x \in \overline{A} \cap
\overline{B}$。第 (3) 条是逆否命题。作为第二个示例，
完整写出第一条分配律的推导：
\[
x \in A \cap (B \cup C)
\iff (x \in A) \land \bigl(x \in B \lor x \in C\bigr)
\iff \bigl(x \in A \land x \in B\bigr) \lor
\bigl(x \in A \land x \in C\bigr),
\]
由\cref{prop:b1:logic:rules}~(6) 的分配律可得，
末一命题即为 $x \in (A \cap B) \cup (A \cap C)$。此类\omterm{def:b1:logic:sets}{集合}恒等式
都可由这一机械的
翻译过程证明——这正是它们无须死记的原因。
\end{proof}

\begin{method}[证明集合相等]\label{met:b1:logic:doubleinc}
要证 $A = B$，只需证明两个包含关系：设 $x \in A$，证明 $x \in
B$；再设 $x \in B$，证明 $x \in A$。此外，也可以串联
等价式 $x \in A \iff \dots \iff x \in B$，
前提是每一步都确为等价。
\end{method}

\begin{omfigure}
\begin{tikzpicture}[scale=0.9]
  \draw[thick] (-2.6,-1.9) rectangle (2.6,1.9);
  \node[below right] at (-2.6,1.9) {$E$};
  \begin{scope}
    \clip (-2.6,-1.9) rectangle (2.6,1.9);
    \fill[omDef!20] (-2.6,-1.9) rectangle (2.6,1.9);
    \fill[white] (-0.7,0) circle (1.1);
    \fill[white] (0.7,0) circle (1.1);
  \end{scope}
  \draw[omDef, thick] (-0.7,0) circle (1.1);
  \draw[omDef, thick] (0.7,0) circle (1.1);
  \node[omDef] at (-1.35,0.5) {$A$};
  \node[omDef] at (1.35,0.5) {$B$};
  \begin{scope}[xshift=6.6cm]
    \draw[thick] (-2.6,-1.9) rectangle (2.6,1.9);
    \node[below right] at (-2.6,1.9) {$E$};
    \begin{scope}
      \clip (-2.6,-1.9) rectangle (2.6,1.9);
      \fill[omThm!20] (-2.6,-1.9) rectangle (2.6,1.9);
      \begin{scope}
        \clip (-0.7,0) circle (1.1);
        \fill[white] (0.7,0) circle (1.1);
      \end{scope}
    \end{scope}
    \draw[omThm, thick] (-0.7,0) circle (1.1);
    \draw[omThm, thick] (0.7,0) circle (1.1);
    \node[omThm] at (-1.35,0.5) {$A$};
    \node[omThm] at (1.35,0.5) {$B$};
  \end{scope}
\end{tikzpicture}

{\small 德摩根律的图示：左图阴影区域为
 $\overline{A \cup B} = \overline A \cap \overline B$（两圆之外的
部分）；右图为 $\overline{A \cap B} = \overline
A \cup \overline B$（除去透镜形重叠部分之外的全部）。
图示本身并非证明，但它使
\cref{prop:b1:logic:setalgebra}中逐元素推演的证明难以记错。}
\end{omfigure}

\section{映射}

\begin{definition}[映射、像、原像]\label{def:b1:logic:map}
称 $f \colon E \to F$ 为\emph{映射}\index{映射}（或\emph{函数}），
若它对每个元素 $x$（取自\omterm{def:b1:logic:sets}{集合} $E$，该\omterm{def:b1:logic:sets}{集合}称为\emph{定义域}），
恰好指定一个元素 $f(x)$，它属于\omterm{def:b1:logic:sets}{集合} $F$（称为\emph{陪域}）。对 $A
\subseteq E$ 与 $B \subseteq F$：
\[
f(A) = \{f(x) : x \in A\} \subseteq F,
\qquad
f^{-1}(B) = \{x \in E : f(x) \in B\} \subseteq E
\]
分别称为 $A$ 的\emph{直像}与 $B$ 的\emph{原像}
\index{原像}。映射 $f \colon E \to F$ 与
$g \colon F \to G$ 的\emph{复合}为 $g \circ f \colon E \to G$，$x \mapsto
g(f(x))$。
\end{definition}

\begin{remark}
记号 $f^{-1}(B)$ \emph{并不}预设逆\omterm{def:b1:logic:map}{映射}的存在：
$f^{-1}(B)$ 对每个 $f$ 都有定义。\omterm{def:b1:logic:map}{原像}的性质比像更好：
$f^{-1}$ 保持并、交与补运算，而
$f(A \cap A') \subseteq f(A) \cap f(A')$ 可能是严格包含
（见\cref{exo:b1:logic:8}）。
\end{remark}

\begin{example}[计算像与原像]\label{ex:b1:logic:images}
设 $f \colon \R \to \R$，$x \mapsto x^2$。则：
\[
f\bigl(\intcc{-1}{2}\bigr) = \intcc04, \qquad
f^{-1}\bigl(\intcc14\bigr) = \intcc{-2}{-1} \cup \intcc12, \qquad
f^{-1}(\{-1\}) = \emptyset .
\]
对第一式：每个 $x \in \intcc{-1}2$ 都有 $x^2 \in \intcc04$，
反之，每个 $y \in \intcc04$ 都可表示为 $y = (\sqrt y)^2$，且 $\sqrt y
\in \intcc02 \subseteq \intcc{-1}2$——注意，像
\emph{并非} $\intcc14 = \{(-1)^2, 2^2\}$，区间的像
不能仅由端点决定。对第二式：$1 \leq x^2 \leq 4 \iff
1 \leq \abs x \leq 2$，这分为两段。第三式
说明\omterm{def:b1:logic:map}{原像}可以为空集——$f^{-1}(B)$ 总是有意义的，
无论 $B$ 与像的交集多么小。最后，
在本例中观察上面注记所提到的严格包含现象：
取 $A = \intcc{-1}0$ 与 $A' = \intcc01$，则 $f(A \cap A') =
f(\{0\}) = \{0\}$，而 $f(A) \cap f(A') = \intcc01$。
\end{example}

\begin{definition}[单射、满射、双射]\label{def:b1:logic:inj}
\omterm{def:b1:logic:map}{映射} $f \colon E \to F$ 具有如下性质：
\begin{itemize}
  \item 不同元素的像互不相同时，为\emph{单射}\index{单射}，即
        像互不相同： $\forall x, x' \in E,\ f(x) = f(x') \implies
        x = x'$；
  \item $F$ 中每个元素都被取到时，为\emph{满射}\index{满射}，即
        $\forall y \in F,\ \exists x \in E,\ f(x) = y$；
  \item 既为单射又为满射时，为\emph{双射}\index{双射}，即
        每个 $y \in F$ 恰有一个\omterm{def:b1:logic:map}{原像}。
\end{itemize}
\end{definition}

\begin{theorem}[逆映射]\label{thm:b1:logic:inverse}
\omterm{def:b1:logic:map}{映射} $f \colon E \to F$ 为\omterm{def:b1:logic:inj}{双射}，当且仅当存在\omterm{def:b1:logic:map}{映射}
$g \colon F \to E$，使得 $g \circ f = \mathrm{id}_E$ 且 $f \circ g =
\mathrm{id}_F$。此时 $g$ 唯一，记作 $f^{-1}$，并
称 $f$ 的\emph{逆\omterm{def:b1:logic:map}{映射}}，且 $f^{-1}$ 本身也是\omterm{def:b1:logic:inj}{双射}，
并有 $(f^{-1})^{-1} = f$。
\end{theorem}

\begin{proof}
($\Rightarrow$) 若 $f$ 为\omterm{def:b1:logic:inj}{双射}，则每个 $y \in F$ 都有唯一的
\omterm{def:b1:logic:map}{原像}；令 $g(y)$ 为该\omterm{def:b1:logic:map}{原像}。由构造，$f(g(y)) = y$，
且 $g(f(x)) = x$，因为 $x$ 正是\emph{那个}
$f(x)$ 的\omterm{def:b1:logic:map}{原像}。

($\Leftarrow$) 假设存在这样的 $g$。若 $f(x) = f(x')$，则对其两边
作用 $g$ 得 $x = x'$，故 $f$ 为\omterm{def:b1:logic:inj}{单射}。对 $y \in F$，取 $x = g(y)$
满足 $f(x) = y$，故 $f$ 为\omterm{def:b1:logic:inj}{满射}。

唯一性：若 $g$ 与 $h$ 均满足要求，则 $g = g \circ \mathrm{id}_F
= g \circ (f \circ h) = (g \circ f) \circ h = h$。最后，这两个恒等式
关于 $f$ 与 $g$ 对称，故 $g = f^{-1}$ 为\omterm{def:b1:logic:inj}{双射}，
且其逆为 $f$。
\end{proof}

\begin{example}[实际计算逆映射]\label{ex:b1:logic:inversecomp}
设 $f \colon \R \to \intoo0{+\infty}$，$f(x) = \eu^{2x+1}$。为求其逆，
解方程 $y = f(x)$，其中 $y > 0$ 为给定值：
\[
y = \eu^{2x+1} \iff \ln y = 2x + 1 \iff x = \frac{\ln y - 1}2 ,
\]
每一步在所声明的定义域上都可逆。
此计算一次性给出全部结论：对陪域中的每个 $y$
恰有一个解 $x$，故 $f$ 为\omterm{def:b1:logic:inj}{双射}，
并有
\[
f^{-1} \colon \intoo0{+\infty} \to \R,
\qquad
f^{-1}(y) = \frac{\ln y - 1}2 .
\]
对两个复合作快速检验（$f^{-1}(f(x)) = \frac{(2x+1) -
1}2 = x$ 与 $f(f^{-1}(y)) = \eu^{\ln y} = y$），
验证\cref{thm:b1:logic:inverse}的判别准则。要点在于：“解出
$x$ 并观察\omterm{def:b1:logic:equiv}{等价关系}”同时就是存在性证明，
唯一性证明与公式——但只有在
陪域被正确声明时才成立（$f$ \emph{不}是
到 $\R$ 上的\omterm{def:b1:logic:inj}{满射}）。
\end{example}

\begin{proposition}[复合与三种性质]\label{prop:b1:logic:comp}
设 $f \colon E \to F$，$g \colon F \to G$。
\begin{enumerate}
  \item 若 $f$ 与 $g$ 均为\omterm{def:b1:logic:inj}{单射}（\omterm{def:b1:logic:inj}{满射}、\omterm{def:b1:logic:inj}{双射}的情形同理），则
        $g \circ f$ 也具有该性质；在\omterm{def:b1:logic:inj}{双射}情形下，$(g \circ f)^{-1} = f^{-1} \circ
        g^{-1}$。
  \item 若 $g \circ f$ 为\omterm{def:b1:logic:inj}{单射}，则 $f$ 为\omterm{def:b1:logic:inj}{单射}。若
        $g \circ f$ 为\omterm{def:b1:logic:inj}{满射}，则 $g$ 为\omterm{def:b1:logic:inj}{满射}。
\end{enumerate}
\end{proposition}

\begin{proof}
(1) 若 $g(f(x)) = g(f(x'))$，由 $g$ 的\omterm{def:b1:logic:inj}{单射}性得 $f(x) = f(x')$，
再由 $f$ 的\omterm{def:b1:logic:inj}{单射}性得 $x = x'$。若 $z \in G$，由
$g$ 的\omterm{def:b1:logic:inj}{满射}性得 $y$，满足 $g(y) = z$，再由 $f$ 的\omterm{def:b1:logic:inj}{满射}性得 $x$，使
$f(x) = y$，于是 $g(f(x)) = z$。在\omterm{def:b1:logic:inj}{双射}情形下，直接验证
$f^{-1} \circ g^{-1}$ 是 $g \circ f$ 的双侧逆，再由
\cref{thm:b1:logic:inverse}中的唯一性即可得证。

(2) 若 $f(x) = f(x')$，则 $g(f(x)) = g(f(x'))$，由
$g \circ f$ 的\omterm{def:b1:logic:inj}{单射}性得 $x = x'$。若 $z \in G$，由 $g \circ f$ 的\omterm{def:b1:logic:inj}{满射}性
得 $x$，使 $g(f(x)) = z$；于是 $y = f(x)$ 满足 $g(y) = z$。
\end{proof}

\begin{example}[(2) 不可加强]\label{ex:b1:logic:compsharp}
在\cref{prop:b1:logic:comp}~(2) 中，不能加强其
结论：$g \circ f$ 为\omterm{def:b1:logic:inj}{双射}，并\emph{不}能推出 $f$
为\omterm{def:b1:logic:inj}{满射}，或 $g$ 为\omterm{def:b1:logic:inj}{单射}。取 $E = G = \{1\}$，$F = \{1,
2\}$，且 $f(1) = 1$，$g(1) = g(2) = 1$，则 $g \circ f =
\mathrm{id}_E$ 为\omterm{def:b1:logic:inj}{双射}，但 $f$ 未取到元素 $2$，并且
$g$ 把两个元素映为同一个值。由此得到一条精确的记账规则：
复合中的信息关于\omterm{def:b1:logic:inj}{单射}性流向\emph{内层}\omterm{def:b1:logic:map}{映射}，
关于\omterm{def:b1:logic:inj}{满射}性则流向\emph{外层}\omterm{def:b1:logic:map}{映射}，绝不反过来。
（\cref{exo:b1:logic:9}用无限集构造了同一现象，
它正是单侧逆背后的动力，
亦即单侧逆的来源。）
\end{example}

\begin{example}\label{ex:b1:logic:injexamples}
$f \colon \R \to \R$，$x \mapsto x^2$ 既不是\omterm{def:b1:logic:inj}{单射}（$f(-1) =
f(1)$），也不是\omterm{def:b1:logic:inj}{满射}（$-1$ 没有\omterm{def:b1:logic:map}{原像}）。限制定义域与
陪域后，$f \colon \R_+ \to \R_+$，$x \mapsto x^2$ 为\omterm{def:b1:logic:inj}{双射}，
其逆为 $y \mapsto \sqrt y$。\omterm{def:b1:logic:map}{映射}的\omterm{def:b1:logic:inj}{单射}性或\omterm{def:b1:logic:inj}{满射}性
取决于所声明的定义域与陪域，而不仅仅取决于公式。
\end{example}

\section{关系}

\begin{definition}[等价关系]\label{def:b1:logic:equiv}
设\emph{二元关系} $\mathcal{R}$ 是\omterm{def:b1:logic:sets}{集合} $E$ 上的
\emph{等价关系}\index{等价关系}，当且仅当它满足：
\emph{自反的}（$x \mathbin{\mathcal{R}} x$，对所有 $x$）、
\emph{对称的}（$x \mathbin{\mathcal{R}} y \implies y
\mathbin{\mathcal{R}} x$），且\emph{传递的}（$x
\mathbin{\mathcal{R}} y$ 且 $y \mathbin{\mathcal{R}} z$ 蕴含 $x
\mathbin{\mathcal{R}} z$）。$x$ 的\emph{等价类}为
$\mathrm{cl}(x) = \{y \in E : x \mathbin{\mathcal{R}} y\}$。
\end{definition}

\begin{example}[检验三条公理]\label{ex:b1:logic:fracpart}
在 $\R$ 上，规定 $x \mathbin{\mathcal{R}} y$，当 $x - y \in \Z$ 时。
\emph{自反性：}$x - x = 0 \in \Z$。\emph{对称性：}若 $x - y
\in \Z$，则 $y - x = -(x - y) \in \Z$。\emph{传递性：}若 $x -
y \in \Z$ 且 $y - z \in \Z$，则 $x - z = (x - y) + (y - z) \in
\Z$（两个整数之和）。因此 $\mathcal R$ 是一个\omterm{def:b1:logic:equiv}{等价
关系}，且 $\mathrm{cl}(x) = x + \Z = \{x + k : k \in \Z\}$：
每个\omterm{def:b1:logic:equiv}{等价类}恰含 $\intco01$ 中的一个代表元，称其为
\emph{小数部分}。与之相对，关系“$\abs{x - y}
\leq 1$”在 $\R$ 上具有自反性与对称性，但\emph{不}
具有传递性（$0 \mathbin{\mathcal R} 1$ 且 $1 \mathbin{\mathcal R}
2$，而 $\abs{0 - 2} > 1$）：接近性不会传递，也不存在
\omterm{def:b1:logic:equiv}{等价类}的\omterm{thm:b1:logic:partition}{划分}——这是一个有用的反例，值得在
验证公理显得例行公事时牢记。
\end{example}

\begin{theorem}[等价类构成划分]\label{thm:b1:logic:partition}
设 $\mathcal{R}$ 为 $E$ 上的\omterm{def:b1:logic:equiv}{等价关系}。则
\omterm{def:b1:logic:equiv}{等价类}均非空，两两不交或相等，且它们的
并为 $E$：它们构成 $E$ 的一个\emph{划分}\index{划分}。
反之，$E$ 的每个\omterm{thm:b1:logic:partition}{划分}恰好由一个\omterm{def:b1:logic:equiv}{等价关系}这样产生
（“位于同一块中”）。
\end{theorem}

\begin{proof}
由自反性，$x \in \mathrm{cl}(x)$，故\omterm{def:b1:logic:equiv}{等价类}非空，且
其并为 $E$。设 $\mathrm{cl}(x) \cap \mathrm{cl}(y) \neq \emptyset$，
不妨设 $z$ 同属两者，则 $x \mathbin{\mathcal{R}} z$ 且 $y
\mathbin{\mathcal{R}} z$，由对称性与传递性得 $x
\mathbin{\mathcal{R}} y$。现任取 $t \in \mathrm{cl}(y)$，
由传递性得 $t \in \mathrm{cl}(x)$，对称地，两类
相等。下证逆命题：设 $(E_i)_{i \in I}$ 为 $E$ 的一个
\omterm{thm:b1:logic:partition}{划分}，并定义 $x \mathbin{\mathcal S} y$ 表示
“某一块同时含有 $x$ 与 $y$”。\emph{自反性：}$x$ 属于某块，
该块便两次含有 $x$。
\emph{对称性：}定义条件关于 $x$ 与
$y$ 对称。\emph{传递性：}若 $x, y \in E_i$ 且 $y, z \in E_j$，则
$y \in E_i \cap E_j$，故 $E_i = E_j$（不同的块互不相交），
从而 $x, z$ 同属一块。在 $\mathcal S$ 下，$x$ 的\omterm{def:b1:logic:equiv}{等价类}
恰为含 $x$ 的那一块，因此\omterm{def:b1:logic:equiv}{等价类}即为给定的
各块。最后，关系由其\omterm{def:b1:logic:equiv}{等价类}决定：两个
具有相同\omterm{def:b1:logic:equiv}{等价类}的\omterm{def:b1:logic:equiv}{等价关系}所联系的数对相同，
因为二者各自使 $x$ 与 $y$ 相关，当且仅当 $y$ 属于
$x$ 所在的类——由此即得唯一性。
\end{proof}

\begin{example}\label{ex:b1:logic:congruence}
在 $\Z$ 上，模 $n$ 同余（$x \equiv y \pmod n$，当 $n$ 整除
$x - y$）是一个\omterm{def:b1:logic:equiv}{等价关系}；其\omterm{def:b1:logic:equiv}{等价类}即为 $n$ 个
整数集，同一集中的整数被 $n$ 除余数相同。本例
在\cref{ch:b1:structures}中将成为环 $\Z/n\Z$。
\end{example}

\begin{definition}[序关系]\label{def:b1:logic:order}
关系 $\preceq$ 是 $E$ 上的\emph{序关系}\index{序关系}，当它
满足自反性、\emph{反对称性}（$x \preceq y$ 且 $y
\preceq x$ 蕴含 $x = y$），且满足传递性。若任意两元素
均可比较则称该序为\emph{全序}，否则为\emph{偏序}。元素
$M \in A \subseteq E$ 称为 $A$ 的\emph{最大元}，当
$a \preceq M$ 对所有 $a \in A$ 成立时；最大元（及最小元）
若存在，则必唯一。
\end{definition}

\begin{example}\label{ex:b1:logic:orders}
$(\R, \leq)$ 是全序集。$(\mathcal{P}(E), \subseteq)$ 当
$E$ 有两个元素时即为偏序集：$\{a\}$ 与
$\{b\}$ 不可比较。记 $A = \{\{a\}, \{b\}\}$，它是
$\mathcal{P}(\{a,b\})$ 的子集，无最大元，却有一个上界
$\{a, b\}$。最大元与上界之间的区别
对 $\R$ 而言，这一区别在\cref{ch:b1:reals}中再次出现。
\end{example}

\begin{example}[网格 $\N^2$ 上的两种序]\label{ex:b1:logic:lexorder}
在自然数对上逐分量比较：$(a, b) \preceq (a',
b')$，当 $a \leq a'$ \emph{且} $b \leq b'$ 时（称为\emph{积序}
）。这是一个序——每条公理都逐分量继承，
但它只是偏序：$(1, 3)$ 与 $(2, 0)$
不可比较。现在像查字典那样比较：$(a, b)
\preceq_{\mathrm{lex}} (a', b')$，当 $a < a'$，或 $a = a'$ 且 $b
\leq b'$ 时（称为\emph{字典序}）。传递性需要
分两种情形验证，但它成立，并且任意两对数现在都可比较：
该序为全序。这两种序对同一\omterm{def:b1:logic:sets}{集合}给出的排序不同
——$(0, 100) \preceq_{\mathrm{lex}} (1, 0)$，而积序对此未作判断，
——这提醒我们：序是一种人所\emph{选取}的结构，
而非\omterm{def:b1:logic:sets}{集合}本身的性质。字典序
比较，也是把若干排序标准
合并为一个的标准技巧。
\end{example}

\begin{remark}[插曲：以双射刻画大小]\label{rem:b1:logic:interlude}
本章一个不太显眼的主题值得特别指出：\omterm{def:b1:logic:inj}{双射}
正是数学家对“大小相同”的刻画。对有限集而言，这一点
便化为\cref{ch:b1:counting}中的计数演算，其中
每个公式实质上都是一个\omterm{def:b1:logic:inj}{双射}；对无限集而言，它
则化为下文的周末问题，其中 $\N$、$\Q$ 与 $\R$ 呈现
出真正不同的大小。同一对应在本卷中还将以更精细的形式
再出现两次：数列
（见\cref{ch:b1:seq}）实质上只是\omterm{def:b1:logic:map}{映射} $\N \to \R$，因而关于数列的命题
就是关于一族\omterm{def:b1:logic:map}{映射}的命题；线性代数则
不以\omterm{def:b1:logic:inj}{双射}而以\emph{线性}\omterm{def:b1:logic:inj}{双射}来衡量向量空间，
其存在性由一个数决定，
即维数（见\cref{ch:b1:findim}）。每当遇到新的
“同一性”——如等势、群的同构
（见\cref{ch:b1:structures}）、线性同构——时，
\cref{thm:b1:logic:inverse}中的模式便会重复：“同一性”是一个
可逆的、保持结构的\omterm{def:b1:logic:map}{映射}。
\end{remark}

\begin{remark}[本章的应用]\label{rem:b1:logic:whereused}
处处都用到——但有几处值得特别指出。三重量词
的技巧（见\cref{ex:b1:logic:limit}）是
\cref{ch:b1:seq,ch:b1:continuity}的日常工具：每个极限证明都是
针对任意 $\varepsilon$ 进行的一场博弈。\omterm{def:b1:logic:equiv}{等价类}以
$\Z/n\Z$ 中的同余类的形式再次出现，见\cref{ch:b1:structures}，其中
\cref{thm:b1:logic:partition}中的\omterm{thm:b1:logic:partition}{划分}获得了自身的代数
结构。\omterm{def:b1:logic:order}{序关系}、上界与上确界
成为 $\R$ 的公理化核心，在\cref{ch:b1:reals}中展开。
\omterm{def:b1:logic:inj}{单射}、\omterm{def:b1:logic:inj}{满射}与\omterm{def:b1:logic:inj}{双射}将再次出现，
即\cref{ch:b1:linmaps}中的线性\omterm{def:b1:logic:map}{映射}，其中\omterm{def:b1:logic:inj}{单射}性可在单个
向量（即核）上检验；而下文的周末问题将朴素的
\omterm{def:b1:logic:inj}{双射}概念发展为一门关于\emph{无限集的
大小}的理论，其结论（$\Q$ 的可数性、
$\R$ 的不可数性）在\cref{ch:b1:reals,ch:b1:topology}中重新出现。
\end{remark}

\section{练习}

\begin{exercise}[$\star$]\label{exo:b1:logic:1}
写出下列每个命题的否定，且不使用“不”“非”等否定词：
\begin{enumerate}
  \item $\forall x \in \R,\ \exists y \in \R,\ x + y > 0$；
  \item $\exists x \in \R,\ \forall y \in \R,\ xy = 0$；
  \item $\forall \varepsilon > 0,\ \exists \delta > 0,\ \forall x \in
        \R,\ \abs{x} \leq \delta \implies \abs{f(x)} \leq \varepsilon$
        （对固定\omterm{def:b1:logic:map}{映射} $f \colon \R \to \R$）。
\end{enumerate}
然后判断命题 (1) 与 (2) 是否为真。
\end{exercise}

\begin{exercise}[$\star$]\label{exo:b1:logic:2}
设 $P, Q$ 为命题。利用真值表证明
$\lnot(P \implies Q) \iff P \land (\lnot Q)$，并由此推出
下述命题的否定：“若函数可导，则它连续”。
\end{exercise}

\begin{exercise}[$\star$]\label{exo:b1:logic:3}
用逆否命题证明：对 $x \in \R$，若 $x^3 + x \geq 2$，则
$x \geq 1$。再用反证法证明：不存在最小的
严格正实数。
\end{exercise}

\begin{exercise}[$\star$]\label{exo:b1:logic:4}
用数学归纳法证明，对所有 $n \in \N$：
\begin{enumerate}
  \item $\sum_{k=0}^{n} 2^k = 2^{n+1} - 1$；
  \item $4^n + 5$ 能被 $3$ 整除。
\end{enumerate}
\end{exercise}

\begin{exercise}[$\star$]\label{exo:b1:logic:5}
找出下述“证明”中的错误，它要证明所有铅笔颜色相同。
\emph{设 $P(n)$：“在任意 $n$ 支铅笔的\omterm{def:b1:logic:sets}{集合}中，所有铅笔
颜色相同”。$P(1)$ 显然成立。假设 $P(n)$ 成立，取 $n+1$
支铅笔；去掉最后一支，前 $n$ 支颜色相同；
去掉第一支，后 $n$ 支颜色相同；于是
所有 $n+1$ 支铅笔颜色相同。}
\end{exercise}

\begin{exercise}[$\star$]\label{exo:b1:logic:6}
设 $A, B, C$ 为 $E$ 的子集。证明：
\begin{enumerate}
  \item $A \setminus B = A \cap \overline{B}$；
  \item $(A \cup B) \setminus C = (A \setminus C) \cup (B \setminus
        C)$；
  \item $A \subseteq B \iff A \cup B = B \iff A \cap B = A$。
\end{enumerate}
\end{exercise}

\begin{exercise}[$\star\star$]\label{exo:b1:logic:7}
对下列每个\omterm{def:b1:logic:map}{映射}，判断（并证明）它是否为\omterm{def:b1:logic:inj}{单射}、\omterm{def:b1:logic:inj}{满射}、
\omterm{def:b1:logic:inj}{双射}：
\begin{enumerate}
  \item $f \colon \N \to \N$，$n \mapsto n + 1$；
  \item $g \colon \Z \to \Z$，$n \mapsto n + 1$；
  \item $h \colon \R \setminus \{1\} \to \R$，$x \mapsto
        \frac{x+1}{x-1}$。
\end{enumerate}
对 $h$，调整陪域使其成为\omterm{def:b1:logic:inj}{双射}，并求
逆\omterm{def:b1:logic:map}{映射}。
\end{exercise}

\begin{exercise}[$\star\star$]\label{exo:b1:logic:8}
设 $f \colon E \to F$，又设 $A, A' \subseteq E$ 及 $B, B' \subseteq
F$。
\begin{enumerate}
  \item 证明 $f^{-1}(B \cap B') = f^{-1}(B) \cap f^{-1}(B')$ 及
        $f(A \cup A') = f(A) \cup f(A')$。
  \item 证明 $f(A \cap A') \subseteq f(A) \cap f(A')$，并给出一个
        使包含关系严格成立的例子。
  \item 证明：$f$ 为\omterm{def:b1:logic:inj}{单射}，当且仅当 $f(A \cap A') = f(A)
        \cap f(A')$ 对所有 $A, A'$ 成立。
\end{enumerate}
\end{exercise}

\begin{exercise}[$\star\star$]\label{exo:b1:logic:9}
设 $f \colon E \to F$ 与 $g \colon F \to E$ 满足 $g \circ f =
\mathrm{id}_E$。证明 $f$ 为\omterm{def:b1:logic:inj}{单射}，$g$ 为\omterm{def:b1:logic:inj}{满射}。
给出一个例子，使 $f$ 与 $g$ 均不是\omterm{def:b1:logic:inj}{双射}。
\end{exercise}

\begin{exercise}[$\star\star$]\label{exo:b1:logic:10}
在 $\R$ 上定义 $x \mathbin{\mathcal{R}} y \iff x^2 - y^2 = x - y$。
证明 $\mathcal{R}$ 为\omterm{def:b1:logic:equiv}{等价关系}，并描述
每个实数 $x$ 的\omterm{def:b1:logic:equiv}{等价类}。哪些\omterm{def:b1:logic:equiv}{等价类}恰含一个
元素？
\end{exercise}

\begin{exercise}[$\star\star\star$]\label{exo:b1:logic:11}
（康托尔）设 $E$ 为\omterm{def:b1:logic:sets}{集合}，证明不存在从 $E$
到 $\mathcal{P}(E)$ 的\omterm{def:b1:logic:inj}{满射}。\emph{提示：给定 $f \colon E \to
\mathcal{P}(E)$，考虑 $D = \{x \in E : x \notin f(x)\}$。}
\end{exercise}

\begin{exercise}[$\star\star\star$]\label{exo:b1:logic:12}
设 $f \colon E \to F$ 为\omterm{def:b1:logic:map}{映射}。定义 $\Phi \colon \mathcal{P}(F)
\to \mathcal{P}(E)$，使 $\Phi(B) = f^{-1}(B)$。
\begin{enumerate}
  \item 证明：$f$ 为\omterm{def:b1:logic:inj}{满射}，当且仅当 $\Phi$
        为\omterm{def:b1:logic:inj}{单射}。
  \item 证明：$f$ 为\omterm{def:b1:logic:inj}{单射}，当且仅当 $\Phi$
        为\omterm{def:b1:logic:inj}{满射}。
\end{enumerate}
\end{exercise}

\section{周末问题：比较无穷}

\begin{problem}\label{pb:b1:logic:1}
两个\omterm{def:b1:logic:sets}{集合}何时才算“元素个数相同”？康托尔的
回答——即二者之间存在\omterm{def:b1:logic:inj}{双射}——竟也
适用于无穷\omterm{def:b1:logic:sets}{集合}，并把无穷\omterm{thm:b1:logic:partition}{划分}为本质不同的“大小”。本题从
本章的基本定义出发，建立起全部工具：
康托尔--施罗德--伯恩斯坦定理（两个\omterm{def:b1:logic:inj}{单射}可构造出一个\omterm{def:b1:logic:inj}{双射}）、
$\Q$ 的可数性、借助对角线法证明 $\R$ 不可数，以及
康托尔在 1874 年的惊人结论：
\emph{\omterm{pb:b1:logic:1}{超越数}确实存在，
而且数量多得惊人}，却无需写出哪怕一个具体的\omterm{pb:b1:logic:1}{超越数}。
\index{可数集}\index{康托尔--施罗德--伯恩斯坦定理}
下文中，对\omterm{def:b1:logic:sets}{集合} $E$ 与 $F$，记 $E \preceq F$，若存在
从 $E$ 到 $F$ 的\omterm{def:b1:logic:inj}{单射}；并记 $E \approx F$（“$E$ 与
$F$ \emph{等势}”），若存在从 $E$
到 $F$ 的\omterm{def:b1:logic:inj}{双射}。

\medskip
\noindent\textbf{第一部分——比较的词汇。}
\begin{enumerate}
  \item 证明 $\approx$ 具有\omterm{def:b1:logic:equiv}{等价关系}的性质：
        $E \approx E$；若 $E \approx F$，则 $F \approx E$；若
        $E \approx F$ 且 $F \approx G$，则 $E \approx G$。（请准确引用
        \cref{thm:b1:logic:inverse}与
        \cref{prop:b1:logic:comp}。）
  \item 证明 $\preceq$ 具有传递性，并证明
        \omterm{def:b1:logic:inj}{单射} $f \colon E \to F$ 总满足 $E \approx f(E)$。
  \item 设 $E \neq \emptyset$。证明：$E \preceq F$ 当且仅当
        存在从 $F$ 到 $E$ 的\omterm{def:b1:logic:inj}{满射}。
  \item 验证 $n \mapsto n + 1$ 是从 $\N$ 到
        $\N^* = \N \setminus \{0\}$ 上的\omterm{def:b1:logic:inj}{双射}，并验证
        \[
        \sigma(n) = \frac n2 \ \ (n \text{ 为偶数}), \qquad
        \sigma(n) = -\frac{n+1}2 \ \ (n \text{ 为奇数})
        \]
        是从 $\N$ 到 $\Z$ 的\omterm{def:b1:logic:inj}{双射}。由此可见，去掉一个点，
        或将其扩充为包含负整数的\omterm{def:b1:logic:sets}{集合}，都不改变
        $\N$ 的大小。
\end{enumerate}

\medskip
\noindent\textbf{第二部分——康托尔--施罗德--伯恩斯坦
定理。} 设 $f \colon E \to F$ 与 $g \colon F \to E$ 为两个
\omterm{def:b1:logic:inj}{单射}。定义
\[
C_0 = E \setminus g(F), \qquad C_{n+1} = g\bigl(f(C_n)\bigr)
\ \ (n \in \N), \qquad C = \bigcup_{n \in \N} C_n,
\]
并令 $h \colon E \to F$ 将 $x \in C$ 映为 $f(x)$，将 $x \notin C$
映为唯一的 $y \in F$，使 $g(y) = x$。
\begin{enumerate}[resume]
  \item 验证 $h$ 是良定义的：若 $x \notin C$，则 $x \in
        g(F)$，且元素 $y$ 使 $g(y) = x$ 成立，它是唯一的。
  \item 证明 $g\bigl(f(C)\bigr) = \bigcup_{n \geq 1} C_n
        \subseteq C$。（直像与并集可交换，
        见\cref{exo:b1:logic:8}。）
  \item 证明 $h$ 是\omterm{def:b1:logic:inj}{单射}。（分三种情形讨论；在混合情形
        $x \in C$、$x' \notin C$ 下，需说明 $h(x) = h(x')$
        将迫使 $x' \in g(f(C)) \subseteq C$。）
  \item 证明 $h$ 是\omterm{def:b1:logic:inj}{满射}：任给 $y \in F$，区分
        两种情形：$g(y) \notin C$，以及 $g(y) \in C_n$（某 $n
        \geq 1$），并说明 $g(y) \in C_0$ 为何不可能，
        并在每种情形下给出 $y$ 的\omterm{def:b1:logic:map}{原像}。
  \item 以\emph{康托尔--施罗德--伯恩斯坦
        定理}作结：若 $E \preceq F$ 且 $F \preceq E$，则 $E
        \approx F$。请用一句话说明
        此结论为何并非平凡。
  \item 两个应用：（a）证明 $\intcc01 \approx \intoo01$。
        （b）证明 $\varphi(p, q) = 2^p(2q + 1) - 1$ 定义了一个
        从 $\N \times \N$ 到 $\N$ 的\omterm{def:b1:logic:inj}{双射}——\omterm{def:b1:logic:inj}{单射}性由奇偶性论证，
        \omterm{def:b1:logic:inj}{满射}性用强归纳法证明，
        （见\cref{thm:b1:logic:induction}）。由此 $\N \times \N \approx
        \N$：整点构成的平面并不比
        直线更大。
\end{enumerate}

\medskip
\noindent\textbf{第三部分——可数集。}称\omterm{def:b1:logic:sets}{集合} $E$
为\emph{至多可数}，当 $E \preceq \N$ 时；并称
\emph{可数}，当 $E \approx \N$ 时。
\begin{enumerate}[resume]
  \item 证明：每个无穷子集 $A \subseteq \N$
        可数。（递归地定义 $\varphi(n)$ 为
        $A \setminus \{\varphi(0), \dots,
        \varphi(n-1)\}$ 中的最小元；证明 $\varphi$ 严格递增，
        满足 $\varphi(n) \geq n$，并取遍
        $A$ 中的每个元素。）
  \item 由此推得：\omterm{def:b1:logic:sets}{集合}至多可数，当且仅当它是
        有限集或可数集，并注意到第 9 题给出了
        如下捷径：若 $E \preceq \N$ 且 $\N \preceq E$，则 $E$
        可数。
  \item 证明：若 $E$ 与 $F$ 均至多可数，则 $E
        \times F$ 亦至多可数。由此推出 $\Z \times \N^*$ 可数。
  \item 证明 $\Q$ 可数。（将 $\Q$ \omterm{def:b1:logic:inj}{单射}入 $\Z \times
        \N^*$ 中，即将每个有理数写成分母为正的既约分数，
        该表示的唯一性在
        \cref{ch:b1:arith}中证明；再应用第 12 题。）
  \item 证明：可数个至多可数集的并仍是
        至多可数的：若每个 $E_n$（$n \in \N$）均至多可数，则
        $\bigcup_{n \in \N} E_n$ 亦至多可数。（将 $x$ 映为
        数对 $(n, f_n(x))$，其中 $n$ 为\emph{最小}指标，满足
        $x \in E_n$。）
  \item 证明：\emph{有限}子集全体（$\N$ 的子集）构成的\omterm{def:b1:logic:sets}{集合}
        可数。（将有限子集 $F$ 映为 $\sum_{i \in F} 2^i$；
        通过比较两有限集相异处的最大元素来证明\omterm{def:b1:logic:inj}{单射}性，
        其中用到 $\sum_{k=0}^{m-1} 2^k = 2^m - 1$，
        该式见\cref{exo:b1:logic:4}。）
\end{enumerate}

\medskip
\noindent\textbf{第四部分——对角化。}设 $\{0,1\}^{\N}$ 表示
所有\omterm{def:b1:logic:map}{映射} $u \colon \N \to \{0, 1\}$ 构成的\omterm{def:b1:logic:sets}{集合}，即
全体 0-1 序列。
\begin{enumerate}[resume]
  \item 构造 $\mathcal{P}(\N)$ 与
        $\{0,1\}^{\N}$ 之间的\omterm{def:b1:logic:inj}{双射}（用指示函数）。
  \item （对角线论证）设 $\Phi \colon \N \to
        \{0,1\}^{\N}$ 为任一\omterm{def:b1:logic:map}{映射}。考虑数列 $d$，
        其定义为 $d(n) = 1 - \Phi(n)(n)$。证明 $d$ 不在
        $\Phi$ 的像中，并由此断言 $\{0,1\}^{\N}$
        \emph{不}至多可数。请用一句话说明，
        结合第 17 题，这恰是康托尔定理
        （见\cref{exo:b1:logic:11}）在 $E = \N$ 时的特例。
  \item 在本阶段不加证明地承认（中学已熟悉，并在
        \cref{ch:b1:reals}中严格证明）：每个 $x \in
        \intco01$ 都有唯一的\emph{正规}十进小数展开 $x =
        0.d_1 d_2 d_3\dots$（即不以无穷多个
        $9$ 结尾者）。给定任一数列 $(x_n)_{n \geq 1}$，其项取自
        $\intco01$，构造 $x \in \intco01$，使 $x \neq x_n$ 对
        一切 $n$ 成立：取其第 $n$ 位为 $5$，若第 $n$ 位
        数字（取自 $x_n$）不等于 $5$，否则取 $6$。需仔细论证
        $x$ 为正规展开，且避开每个 $x_n$，并由此推得
        $\intco01$ 不是至多可数的。
  \item 由此推得 $\R$ 不可数，且无理数全体构成的\omterm{def:b1:logic:sets}{集合} $\R
        \setminus \Q$ 亦不可数。
        在何种确切意义上，“大多数”实数是无理数？
\end{enumerate}

\medskip
\noindent\textbf{第五部分——康托 1874 年定理：
\omterm{pb:b1:logic:1}{超越数}存在。} 实数 $x$ 称为\emph{代数数}
\index{代数数}，若 $P(x) = 0$ 对某个非零整系数多项式 $P$ 成立，
否则称为\emph{超越数}
\index{超越数}。本部分不加证明地承认以下结论，
其证明见\cref{ch:b1:poly}：次数为 $n$ 的非零多项式
至多有 $n$ 个实根。
\begin{enumerate}[resume]
  \item 证明每个有理数都是\omterm{pb:b1:logic:1}{代数数}，并求出具体的整系数多项式，
        分别以 $\sqrt 2$ 与
        $\sqrt 2 + \sqrt 3$ 为根。
  \item 对固定的 $n \in \N$，证明次数至多为 $n$、系数为整数的多项式全体
        构成的\omterm{def:b1:logic:sets}{集合}可数。
        （将其\omterm{def:b1:logic:inj}{单射}到 $\Z^{n+1}$ 中，并对 $n$ 用归纳法，
        借助第 13 题。）
  \item 由此可知，\emph{全部}整系数多项式构成的\omterm{def:b1:logic:sets}{集合}
        是可数的。
  \item 证明\emph{康托关于\omterm{pb:b1:logic:1}{代数数}的定理}：
        代数实数全体构成的\omterm{def:b1:logic:sets}{集合} $\mathcal{A}$ 可数。
  \item 由此得出：超越实数存在，且\omterm{pb:b1:logic:1}{超越数}全体构成的\omterm{def:b1:logic:sets}{集合}
        不可数。随后用几句话回顾整道题：
        链 $\N \approx
        \Z \approx \Q \approx \mathcal{A}$，严格跃升到 $\R
        \approx$（本质上）$\mathcal{P}(\N)$，其中每个工具
        （康托–施罗德–伯恩斯坦定理、可数并、
        对角线法）各自都起到了决定性作用，并且证明
        存在不可数个\omterm{pb:b1:logic:1}{超越数}，却不点名其中任何一个，这一哲学意味尤为深刻。
        （证明一个\emph{具体}的数，
        例如 $\pi$ 为\omterm{pb:b1:logic:1}{超越数}，要困难得多，
        超出本册的范围。）
\end{enumerate}
\end{problem}
